\documentstyle[%


righttag%


,11pt%


,amssymb%


,russian%               remove in Eng.


,izvru3%                 Set izven% in Eng.


]{amsart}





% remove the dot in STY-file


% Look in hints for \hline in matrix


% 0,5 etc. in Eng.


\newcommand{\tg}{\operatorname{tg}}


\newcommand{\arctg}{\operatorname{arctg}}


\newcommand{\const}{\operatorname{const}}


\newcommand{\diag}{\operatorname{diag}}


\newcommand{\tts}[1]{}





\theoremstyle{plain}


\theorembodyfont{\it}


\newtheorem{assrt}{\hskip\parindent Утверждение}





%\hoffset=-15mm%                  For Eng.


\setcounter{page}{26}


\god{2000}


\nomer{11~(462)} %                        Remove in Eng.


%\nomer{Vol.\,44, No.\,11}%               Set in Eng.


%\str{\pageref{begin}--\pageref{end}}%  Set in Eng.


\udk{519.63}





\author{\it А.В.\,Гулин, А.В.\,Шередина}


% NB:   ^^ remove in Eng.


\title{Границы устойчивости разностных схем}


\thanks{Работа выполнена при финансовой поддержке Российского фонда


фундаментальных исследований (код проекта 99-01-00958).}





\date{}


%\pagestyle{myheadings}%         Set in Eng.


\pagestyle{plain} %              Remove in Eng.


\begin{document}


%\thispagestyle{plain}%          Set in Eng.


\maketitle


\label{begin}





{\bf 1.} {\it Введение.}


Рассматриваются двуслойные и трехслойные операторно-разностные схемы


\begin{gather}


\label{1}


By_{t} +Ay=0, \\


%


\label{2}


By_{\overset{\circ}{t}} +\tau^{2}R y_{\ov tt}+Ay=0,


\end{gather}


где $y=y_{n}=y(t_{n})\in H$


--- функции дискретного аргумента $t_{n}=n\tau$


со значениями в конечномерном линейном пространстве $H$,


$n=0,1,\dotsc$, $\tau >0$, и $A$, $B$, $R$


--- линейные операторы, действующие в $H$.


Здесь обозначено


$$


y_{t}=\frac{y_{n+1}-y_{n}}{\tau},\quad y_{\ov t}=\frac{y_{n}-y_{n-1}}{\tau},


\quad y_{\overset{\circ}{t}}=\frac{y_{n+1}-y_{n-1}}{2\tau}.


$$


Теория устойчивости разностных схем (\ref{1}), (\ref{2}) изложена в [1], [2].


В [3]--[6] получены критерии устойчивости двуслойных и трехслойных


разностных схем с операторными весовыми множителями, т.\,е. схем вида


\begin{gather}


\label{3}


y_{t} +\sigma Ay_{n+1}+(E-\sigma) Ay_{n}=0, \\


%


\label{4}


y_{\ov tt} +\sigma Ay_{n+1}+(E-2\sigma) Ay_{n} + \sigma Ay_{n-1} =0,


\end{gather}


где $\sigma$ --- оператор в $H$, $E$ --- единичный оператор.


Теория разностных схем с операторными множителями имеет свои особенности


и достаточно подробно изложена в [6].





Условия устойчивости, полученные в [3]--[5], характеризуются тем, что их


невозможно ослабить за счет выбора нормы. Приведем необходимые для


дальнейшего теоремы из [3], [5].


Предполагается, что операторы $A$, $B$, $R$


не зависят от номера слоя $n$.


Устойчивость исследуется в нормах $\|y\|_{D}$ =$\sqrt{(Dy,y)}$,


порожденных самосопряженным положительным оператором $D:H\to H$.


В дальнейшем под {\it устойчивостью в пространстве $H_{D}$}


понимается невозрастание со временем $t=t_{n}$


нормы решения $\|y(t_{n})\|_{D}$ разностной задачи.





\begin{thm}


Пусть $A^{*}=A$, $\sigma^{*}=\sigma$.


Если схема $(\ref{3})$ устойчива в каком-либо пространстве $H_{D}$,


то выполнено операторное неравенство


\begin{equation}


\label{5}


A+\tau A\mu A\ge 0,


\end{equation}


где $\mu=\sigma -0,5E$.


Обратно, если выполнено $(\ref{5})$, то схема $(\ref{3})$


устойчива в $H_{A^{2}}$.


\end{thm}





\begin{thm}


Пусть $A^{*}=A$, $\sigma^{*}=\sigma$.


Если схема $(\ref{4})$ устойчива в каком-либо пространстве $H_{D}$,


то выполнено операторное неравенство


%\begin{equation} \label{6}


$$


A+\tau^{2} A\mu A\ge 0,


%\end{equation}


$$


где $\mu=\sigma -0,25E$.


Обратно, если $A+\tau^{2} A\mu A> 0$,


то схема $(\ref{4})$ устойчива в норме


$$


\|y_{n}\|_{D}=\bigg(\bigg\|\frac{y_{n}+y_{n-1}}{2}\bigg\|_{A^{2}}^{2} +


\bigg\|\frac{y_{n}-y_{n-1}}{\tau}\bigg\|_{A+\tau^{2} A\mu A}^{2}


\bigg)^{1/2}.


$$


\end{thm}





{\bf 2.} {\it Границы устойчивости двумерных разностных схем.}


Сформулированные выше теоремы дают возможность провести численное 


исследование устойчивсти конкретных разностных схем,


аппроксимирующих задачи математической физики.


Так, в [4] был получен критерий устойчивости двуслойных схем 


с переменными весовыми множителями для двумерного уравнения 


теплопроводности. Условия устойчивости формулировались в [4] 


в виде требования неотрицательности всех собственных значений 


так называемой матрицы устойчивости, т.\,е.{} некоторой симметричной 


матрицы, порожденной рассматриваемой разностной задачей. 


Неотрицательность собственных значений проверялась численно 


для каждых заданных отношений


$\gamma_{1}=\tau /h_{1}^{2}$ и $\gamma_{2}=\tau /h_{2}^{2}$,


где $\tau$ --- шаг сетки по времени и $h_{1}$, $h_{2}$


--- шаги сетки по пространственным переменным $x_{1}$


и $x_{2}$ соответственно.





Рассмотрим задачу о численном построении границ устойчивости


разностных схем с переменными весовыми множителями для двумерного уравнения


теплопроводности. В прямоугольной


области $(0<x_1<l_1$, $0<x_2<l_2)$


рассматривается многопараметрическое семейство разностных схем


\begin{equation}


\label{7}


\frac{y_{ij}^{n+1}-y_{ij}^{n}}{\tau}=


\sigma_{ij}\Lambda y_{ij}^{n+1} +


(1-\sigma_{ij})\Lambda y_{ij}^{n},


\end{equation}


$i=1,2,\dotsc ,N_1-1$, $j=1,2,\dotsc ,N_2-1$, $h_{1}N_1=l_1$, $h_{2}N_2=l_2$


с нулевыми граничными условиями.


Здесь $\sigma_{ij}$ --- числовые множители


и $\Lambda$ --- пятиточечный разностный оператор,


$\Lambda$=$\Lambda_1+\Lambda_2$, где


\begin{align*}


(\Lambda_1 y)_{ij}^{n}&= \frac{1}{h_{1}}


\bigg(a_{i+1 j}\frac{y_{i+1j}^{n}-y_{i j}^{n}}{h_{1}} -


a_{i j}\frac{y_{ij}^{n}-y_{i-1j}^{n}}{h_{1}} \bigg), \\


%


(\Lambda_2 y)_{ij}^{n}&= \frac{1}{h_{2}}


\bigg(b_{i j+1}\frac{y_{ij+1}^{n}-y_{ij}^{n}}{h_{2}} -


b_{i j}\frac{y_{ij}^{n}-y_{ij-1}^{n}}{h_{2}} \bigg).


\end{align*}


Предполагается, что


$a_{ij} \ge c_1 > 0$, $b_{ij} \ge c_2 > 0$ для всех $i$, $j$.


Уравнение (\ref{7}) зависит от двух сеточных параметров,


$\gamma_{1}=\tau /h_{1}^{2}$ и $\gamma_{2}=\tau /h_{2}^{2}$,


определяющих наряду с $\sigma=\sigma_{ij}$


устойчивость или неустойчивость разностной схемы.





Пусть набор весовых множителей фиксирован. 


Всюду в дальнейшем предполагается, что весовые 


множители $\sigma_{ij}$ не зависят от $\gamma_1$ 


и $\gamma_2$.Примем следующую терминологию.


Разностную схему (\ref{7}) назовем {\it устойчивой в точке}


$(\gamma_{1},\gamma_{2})$ плоскости $\gamma_{1}O\gamma_{2}$,


если при этих значениях $\gamma_{1}$ и


$\gamma_{2}$ схема устойчива. В дальнейшем будем предполагать,


что $\gamma_{1}\ge 0$, $\gamma_{2}\ge 0$, т.\,е.{} ограничимся рассмотрением


первого квадранта плоскости $\gamma_{1}O\gamma_{2}$.


Назовем {\it областью устойчивости} разностной схемы (\ref{7}) множество всех


точек первого квадранта, в которых схема устойчива, и областью 


неустойчивости --- все остальные точки первого квадранта.


{\it Границей устойчивости } разностной схемы (\ref{7})


называется кривая в квадранте $\gamma_{1}\ge 0$, $\gamma_{2}\ge 0$,


разделяющая области устойчивости и неустойчивости. Схема называется абсолютно


устойчивой, если она устойчива во всех точках $(\gamma_{1},\gamma_{2})$


первого квадранта. Тем самым, для абсолютно устойчивой схемы границы


устойчивости не существует.


Для трехслойной разностной схемы


\begin{equation}


\label{8}


\frac{y_{ij}^{n+1}-2y_{ij}^{n}+y_{ij}^{n-1}}{\tau^{2}}=


\sigma_{ij}\Lambda y_{ij}^{n+1} +


(1-2\sigma_{ij})\Lambda y_{ij}^{n}+


\sigma_{ij}\Lambda y_{ij}^{n-1}


\end{equation}


понятие границы устойчивости вводится


точно так же, как и для двуслойной, только в качестве сеточных параметров


$\gamma_{1}$ и $\gamma_{2}$ выступают отношения


$\gamma_{1}=\tau^{2} /h_{1}^{2}$ и $\gamma_{2}=\tau^{2} /h_{2}^{2}$.


\medskip





{\bf 3.} {\it Численное построение границы устойчивости.}


Будем рассматривать для определенности двуслойную схему (\ref{7}).


В работах [7], [8] были численно построены границы устойчивости 


разностных схем (\ref{7}) с $a_{ij} \equiv 1$, $b_{ij} \equiv 1$ 


и с различными наборами весовых множителей.


Для решения уравнения $F(\gamma_1, \gamma_2) = 0$,


определяющего неявно границу устойчивости, применялся итерационный метод


Стеффенсена. Существенным этапом вычислительного алгоритма явился


предложенный Л.Ф.\,Юхно метод нахождения минимального собственного значения


больших разреженных матриц.





Результаты, изложенные в данной работе, основаны на иной, более общей, 


как нам представляется, методике построения границы


устойчивости, не предполагающей применения итерационных методов.


В основе предлагаемого метода лежит переход от декартовых координат


$(\gamma_{1},\gamma_{2})$ к полярным $(r,\varphi)$


и поиск точки границы устойчивости на лучах


$\varphi =\const$.





Схема (\ref{7}) приводится к каноническому виду (\ref{1}), где $A=-\Lambda$.


Известно (см., напр., [1], с.\,236), что матрица $\tau A$ является 


симметричной и положительно определенной. Согласно теореме 1, 


схема (\ref{7}) устойчива тогда и только тогда, когда неотрицательны 


все собственные значения


{\it матрицы устойчивости} $P=(\tau A)+(\tau A)\mu(\tau A)$.


Искомые параметры $\gamma_{1}$ и $\gamma_{2}$ входят только в матрицу


$\tau A$ и не содержатся в матрице $\mu$. Матрица $\tau A$


имеет вид


\begin{equation}


\label{9}


\tau A = \gamma_{1}A_{1}+\gamma_{2}A_{2},


\end{equation}


где $\gamma_{1}=\tau/h_{1}^{2}$, $\gamma_{2}=\tau/h_{2}^{2}$,


$A_1 =-h_{1}^{2} \Lambda_1$, $A_2 = -h_{2}^{2}\Lambda_2$. 


Вид этих матриц зависит от способа перенумерации двумерных массивов 


в одномерный. Пусть, например, $a_{ij} \equiv 1$, $b_{ij} \equiv 1$,


$N_1=4$, $N_2=3$, и массив весовых множителей имеет вид


$$


\sigma =


\begin{bmatrix}


\sigma_{12} & \sigma_{22} & \sigma_{32} \\


\sigma_{11} & \sigma_{21} & \sigma_{31} \\


\end{bmatrix}


.


$$


Тогда, если проводить перенумерацию по строкам, 


то матрицы $A_1$, $A_2$ и $\mu$ имеют вид


\begin{gather*}


A_{1}=


\begin{bmatrix}


2 & -1 & 0 & \vline & 0 & 0 & 0 \\


-1 & 2 & -1 & \vline & 0 & 0 & 0 \\


0 & -1 & 2 & \vline & 0 & 0 & 0\\


\hline


0 & 0 & 0 & \vline & 2 & -1 & 0\\


0 & 0 & 0 & \vline & -1 & 2 & -1 \\


0 & 0 & 0 & \vline & 0 & -1 & 2


\end{bmatrix}


, \ \;


A_{2}=


\begin{bmatrix}


2 & 0 & 0 & \vline & -1 & 0 & 0\\


0 & 2 & 0 & \vline & 0 & -1 & 0\\


0 & 0 & 2 & \vline & 0 & 0 & -1\\


\hline


-1 & 0 & 0 & \vline & 2 & 0 & 0\\


0 & -1 & 0 & \vline & 0 & 2 & 0\\


0 & 0 & -1 & \vline & 0 & 0 & 2


\end{bmatrix}


,


\\


%


\mu=


\begin{bmatrix}


\mu_{11} & 0 & 0 & \vline & 0 & 0 & 0\\


0 & \mu_{21} & 0 & \vline & 0 & 0 & 0\\


0 & 0 & \mu_{31} & \vline & 0 & 0 & 0\\


\hline


0 & 0 & 0 & \vline & \mu_{12} & 0 & 0\\


0 & 0 & 0 & \vline & 0 & \mu_{22} & 0\\


0 & 0 & 0 & \vline & 0 & 0 & \mu_{32}


\end{bmatrix}


,


\end{gather*}


где $\mu_{ij} = \sigma_{ij}-0,5$. 


Важно отметить, что $A_1$ и $A_2$ не зависят от сеточных 


параметров $\gamma_1$, $\gamma_2$. Матрица $\mu$ является диагональной.





Полагая $\gamma_{1}=r\cos\varphi$, $\gamma_{2}=r\sin\varphi$,


получим из (\ref{9}), что


$\tau A =rA_{\varphi}$,


где


%\begin{equation} \label{10}


$$


A_{\varphi}=A_{1}\cos\varphi + A_{2}\sin\varphi.


%\end{equation}


$$


Таким образом, искомый параметр $r$ входит в матрицу $\tau A$ 


в виде числового множителя, а матрица $A_{\varphi}$ не зависит от $r$.


Заметим, что оператор $A_{\varphi}$ отличается от $\tau A$ только 


положительным множителем и, следовательно, является симметричной 


положительно определенной матрицей.


Матрица устойчивости $P=(\tau A)+(\tau A)\mu(\tau


A)$ представляется в виде $P=rP_{\varphi}$, где


\begin{equation}


\label{11}


P_{\varphi}=A_{\varphi} +


rA_{\varphi}\mu A_{\varphi}.


\end{equation}





Поиск границы устойчивости основан на следующих свойствах матрицы


$P_{\varphi}$.





\begin{assrt}


Разностная схема $(\ref{7})$ устойчива в точке


$\gamma_{1}=r\cos\varphi$, $\gamma_{2}=r\sin\varphi$


тогда и только тогда, когда неотрицательны все собственные значения


матрицы $P_{\varphi}$.


\end{assrt}





Действительно, если $r=0$, то $P=rP_{\varphi}=0$ и согласно теореме 1


схема (\ref{7}) устойчива. Если же $r>0$, то неравенства $P\ge 0$ и


$P_{\varphi}\ge 0$ эквивалентны.








Рассмотрим обобщенную задачу на собственные значения


\begin{equation}


\label{12}


A_{\varphi}\mu A_{\varphi}x = \lambda A_{\varphi}x


\end{equation}


и покажем, что свойство знакоопределенности спектра задачи (\ref{12}) 


не зависит от $\varphi$.





\begin{assrt}


Минимальное собственное значение $\lambda_{\min}$ задачи $(\ref{12})$


неотрицательно тогда и только тогда, когда


$\mu_{ij}\ge 0$ для всех $i$ и $j$.


\end{assrt}





\begin{pf}


Заметим, что спектр задачи (\ref{12}) совпадает со 


спектром матрицы


\begin{equation}


\label{13}


M_{\varphi}=A_{\varphi}^{1/2}\mu A_{\varphi}^{1/2}.


\end{equation}





Далее, неотрицательность минимального собственного значения матрицы


$M_{\varphi}$ эквивалентна выполнению матричного неравенства


$M_{\varphi} \ge 0$, которое вследствие невырожденности матрицы


$A_{\varphi}^{1/2}$ эквивалентно неравенству $\mu\ge 0$.


Поскольку $\mu$ --- диагональная матрица, неравенство $\mu\ge 0$


выполнено тогда и только тогда,


когда $\mu_{ij}\ge 0$ для всех $i$ и $j$.


\end{pf}





\begin{assrt}


Если все $\sigma_{ij}\ge 0,5$, то схема


$(\ref{7})$ абсолютно устойчива.


\end{assrt}





\begin{pf} Пусть $\gamma_{1}$ и $\gamma_{2}$


--- любые положительные числа. Определим угол $\varphi =


\arctg (\gamma_{2} /\gamma_{1})$ и для найденного $\varphi$ рассмотрим


матрицу $M_{\varphi}$, построенную согласно (\ref{13}). Из утверждения 2


следует, что минимальное собственное значение $\lambda_{\min}(M_{\varphi})$


матрицы $M_{\varphi}$ неотрицательно. Но тогда для любого $r>0$ все


собственные значения матрицы $E+rM_{\varphi}$ больше или равны единице.


Поэтому выполнено матричное неравенство $E+rM_{\varphi} \ge 0$,


эквивалентное неравенству $A_{\varphi} + rA_{\varphi}\mu A_{\varphi} \ge 0$.


Согласно утверждению 1 схема (\ref{7}) устойчива в точке


$(\gamma_{1},\gamma_{2})$.


\end{pf}





В следующем утверждении найдены координаты точки границы устойчивости,


лежащей на луче $\gamma_2=\gamma_1\tg \varphi$.





\begin{assrt}


Пусть хотя бы один из весовых множителей $\sigma_{ij}< 0,5$.


Тогда при любом $\varphi$ минимальное собственное значение


$\lambda_{\min}(\varphi)$ задачи $(\ref{12})$


отрицательно. Точка с координатами $\gamma_{1}=r\cos\varphi$,


$\gamma_{2}=r\sin\varphi$


принадлежит границе устойчивости разностной схемы $(\ref{7})$


тогда и только тогда, когда $r=-1/\lambda_{\min}(\varphi)$.


\end{assrt}





\begin{pf}


Отрицательность $\lambda_{\min}(\varphi)$ следует из утверждения 2. 


Число $1+r\lambda_{\min}(\varphi)$ является минимальным собственным 


значением матрицы $E+rM_{\varphi}$. Поэтому, если


$r{<}{-}(\lambda_{\min}(\varphi))^{-1}$, то все собственные значения 


матрицы $E+rM_{\varphi}$ положительны и, следовательно, схема (\ref{7}) 


устойчива. Если же $r > - (\lambda_{\min}(\varphi))^{-1}$, то матрица


$E+rM_{\varphi}$ имеет хотя бы одно отрицательное собственное значение


и, следовательно, схема (\ref{7}) неустойчива. Тем самым, при каждом


$\varphi\in [0,\pi /2 ]$ точка $r = - (\lambda_{\min}(\varphi))^{-1}$


и только она принадлежит границе устойчивости разностной схемы (\ref{7}).


\end{pf}





{\bf Следствие.} Если хотя бы один из весовых множителей


$\sigma_{ij} < 0,5$, то существует граница устойчивости $r= r(\varphi)$,


являющаяся однозначной функцией $\varphi$.


\medskip





Таким образом, для определения границы устойчивости разностной схемы


(\ref{7}) в том случае, когда хотя бы один из весовых множителей меньше 0,5,


достаточно при каждом $\varphi\in [0,\pi/2 ]$ найти минимальное собственное


значение $\lambda_{\min}=\lambda_{\min}(\varphi)$ задачи (\ref{12}) и положить


$$


r=-1/\lambda_{\min}(\varphi),\quad


\gamma_{1}=r\cos\varphi, \quad


\gamma_{2}=r\sin\varphi.


$$


При численном построении по переменной $\varphi$


вводится сетка ${\varphi_{k}}$, и вычисления проводятся только


в точках сетки. Заметим, что вычисления можно проводить независимо


друг от друга на каждом луче $\varphi_{k}$.





В случае трехслойной разностной схемы (\ref{8}) сформулированные выше


утверждения претерпевают незначительные изменения. Надо только


заменить условие $\sigma_{ij}\ge 0,5$ условием $\sigma_{ij}\ge 0,25$,


а условие $\sigma_{ij}< 0,5$ --- условием $\sigma_{ij}< 0,25$.


Кроме того, в случае трехслойной схемы имеем $\gamma_1 = \tau^2/h_1^2$


и $\gamma_2 = \tau^2/h_2^2$.


Вычислительный алгоритм построения границы устойчивости остается


без изменения.


\medskip





{\bf 4.} {\it Опорные точки. Базисная прямая.}


Проведенные расчеты показали, что


граница устойчивости, как правило, мало отличается от


отрезка некоторой прямой, определяемой параметрами $\sigma_{ij}$.


Эта прямая пересекается с границей устойчивости в точках


$(\gamma_{10},0)$ и $(0,\gamma_{20})$. Такую


прямую будем называть {\it базисной прямой}, а точки


$(\gamma_{10},0)$ и $(0,\gamma_{20})$ --- {\it опорными


точками}.





Опорные точки можно было бы построить по описанной выше двумерной


методике. Однако в данном случае алгоритм можно существенно


упростить и свести к определению границ устойчивости одномерных


задач. Рассмотрим, например, как находится точка $(\gamma_{10},0)$.


В этом случае полярный угол $\varphi=0$ и матрицы


$A_{\varphi}=A_{1}$, $A_{\varphi}\mu A_{\varphi}$,


входящие в основное уравнение (\ref{12}), становятся блочно-диагональными.


Например, в случае $a_{ij} \equiv 1$, $b_{ij} \equiv 1$,


$N_1=4$, $N_2=3$ матрица $A_{\varphi}\mu A_{\varphi}$


= $A_{1}\mu A_{1}$ имеет вид $A_{1}\mu A_{1}=\diag [M_1,M_2]$,


где


\begin{align*}


M_{1}&=


\begin{bmatrix}


4\mu_{11}+\mu_{21}&-2(\mu_{11}+\mu_{21})&\mu_{21}\\


-2(\mu_{11}+\mu_{21})&\mu_{11}+4\mu_{21}+\mu_{31}&-2(\mu_{21}+\mu_{31})\\


\mu_{21}&-2(\mu_{21}+\mu_{31})&\mu_{21}+4\mu_{31}\\


\end{bmatrix}


, \\


%


M_{2}&=


\begin{bmatrix}


4\mu_{12}+\mu_{22}&-2(\mu_{12}+\mu_{22})&\mu_{22} \\


-2(\mu_{12}+\mu_{22}) &\mu_{12}+4\mu_{22}+\mu_{32} &-2(\mu_{22}+\mu_{32}) \\


\mu_{22} &-2(\mu_{22}+\mu_{32}) & \mu_{22}+4\mu_{32}


\end{bmatrix}


.


\end{align*}





В общем случае


$A_{1}\mu A_{1}=\diag [M_1,M_2, \dotsc ,M_{N_{2}-1}]$,


где $M_j$ --- симметричные матрицы порядка $N_{1}-1$.


Поэтому спектр задачи (10) представляет собой объединение спектров


$N_{2}-1$ матриц, имеющих порядок $N_{1}-1$.


Соответственно в $j$-й диагональной клетке $P_{j}$ матрицы устойчивости


(9) присутствуют элементы только $j$-й строки


$$


\sigma^{(j)}=(\sigma_{1j},\sigma_{2j},\dotsc ,\sigma_{N_{1}-1j})


$$


матрицы $\sigma$. Тем самым, минимальное собственное значение


матрицы $P_{j}$ определяет границу устойчивости $\gamma_{1}(j)$


одномерной задачи с распределением весовых множителей $\sigma^{(j)}$,


а граница устойчивости двумерной задачи определяется как


$$


\gamma_{10}=\min_{1\le j\le N_{2}-1} \gamma_{1}(j).


$$





Тем самым, минимальное из всех чисел $\gamma_{1}(j)$ соответствует худшему


из одномерных вариантов по строкам (понятия ``худшего'' и ``оптимального''


вариантов введены в [9]). Отсюда следует, что


{\it опорная точка


$(\gamma_{10},0)$ не зависит от перестановок строк матрицы $\sigma$}.


Точно так же {\it опорная точка $(0,\gamma_{20})$


определяется худшим вариантом по столбцу матрицы $\sigma$ и


не зависит от перестановки столбцов}.


\medskip





{\bf 5.} {\it Результаты расчетов.} Построение границ устойчивости


проводилось для схем (6), (7) с $a_{ij} \equiv 1$, $b_{ij} \equiv 1$


и с различными наборами весовых множителей $\sigma$.


Поскольку граница устойчивости мало отличается от отрезка базисной прямой,


удобно характеризовать ее тремя величинами: координатами опорных точек


$\gamma_{10}$, $\gamma_{20}$, и функцией $z(\varphi)$ --- отклонением 


границы устойчивости от базисной прямой


в зависимости от угла $\varphi$. Функция $z(\varphi)$ определяется как


$z(\varphi) = r_A(\varphi) - r_B(\varphi)$, где $r_A(\varphi)$ и


$r_B(\varphi)$ --- полярные радиусы точки $A$ границы устойчивости и точки


$B$ базисной прямой, отвечающих одному и тому же полярному углу $\varphi$.


Различным распределениям $\sigma = ( \sigma_{ij})$ соответствуют


различные значения величин $(\gamma_{10}, \gamma_{20}, z(\varphi))$.


В таблицах 1--5 приведены значения опорных точек $\gamma_{10}$, $\gamma_{20}$


и максимального отклонения от базисной прямой


$z = \max\limits_{0\le k\le 10} z( \varphi(k))$


для различных распределений $\sigma =(\sigma_{ij})$.


Здесь $\varphi(k) = kh$, $h = 0,05 \pi$, $k = 0,1,\dotsc,10$.





В таблице 1 показана зависимость границы устойчивости схемы (\ref{7})


от изменения величины весового множителя в одной точке сетки при


неизменных значениях весовых множителей в других точках. В 


проведенных расчетах $N_1=10$, $N_2=7$


и матрица $\sigma$ имела элементы $\sigma_{ij}=1$ для $(i,j) \ne (3,2)$,


$\sigma_{32}=t$, где параметр $t$ менялся от 0 до 0,5.


В таблице 1 приведены данные, отвечающие различным значениям $t$.


Видно, что с увеличением $t$ значения $\gamma_{10}$ и $\gamma_{20}$


резко возрастают, и граница все более удаляется от базисной прямой.





В таблицe 2 приведены значения $\gamma_{10}$, $\gamma_{20}$ и


$z = \max\limits_{0\le k\le 10} z( \varphi(k))$


для различных распределений $\sigma = ( \sigma_{ij})$,


состоящих из нулей и единиц. Здесь $N_1=5$, $N_2=4$.


Заметим, что увеличение числа точек $N_1$ и $N_2$ не приводит к существенному


изменению границы устойчивости. В таблице 3 представлены данные для


распределения $\sigma =(\sigma_{ij})$, аналогичного распределению


шестого столбца таблицы 2, но с числом точек, вдвое б\'ольшим по каждому


направлению. В скобках указаны уточненные по сравнению с таблицей 2 значения


для $N_1=5$, $N_2=4$. Видно, что значения $\gamma_{10}$ и


$z = \max\limits_{0\le k\le 10} z( \varphi(k))$ почти не изменились.





В таблице 4 приведены значения опорных точек


$\gamma_{10}$, $\gamma_{20}$ и максимального отклонения от базисной прямой


$z = \max\limits_{0\le k\le 10} z( \varphi(k))$


для разностных схем, близких к абсолютно устойчивым схемам. Для таких схем


искривление границ устойчивости становится более выраженным, а отклонение от


базисной прямой является величиной $O(1)$.





Во всех вариантах, представленных выше, граница устойчивости была расположена


по одну сторону от базисной прямой, что выражалось в сохранении знака функции


$z(\varphi)$. Вариант, представленный в таблице 5, свидетельствует о том,


что функция $z(\varphi)$ может менять знак. Значения функции $z(\varphi)$


в точках сетки приведены в таблице 6.





Приведем результаты расчетов границ устойчивости трехслойных разностных схем


(7) с различным распределением весовых множителей $\sigma$.


Известно, что явная схема ($\sigma_{ij}=0$ для всех $i$ и $j$) устойчива


при условии $\gamma_1 a_1 + \gamma_2 a_2 < 1$, где $\gamma_1 = \tau^2/h_1^2$,


$\gamma_2 = \tau^2/h_2^2$ и $a_1 = \cos^2 \frac{\pi}{2N_1}$,


$a_2 = \cos^2 \frac{\pi}{2N_2}$. Граница устойчивости представляет собой


прямую $\gamma_1 a_1 + \gamma_2 a_2 = 1$. Результаты расчетов представлены


во второй строке таблицы 7. Сравнение с точным значением показывает,


что максимальная погрешность в определении границы равна $1,7\cdot 10^{-8}$.


В третьей строке таблицы 7 представлены данные о границе устойчивости


трехслойной схемы (\ref{8}) с шахматным распределением весовых множителей


1 и 0,24. Такая схема близка к абсолютно устойчивой схеме


(абсолютная устойчивость наступает при замене параметра 0,24 на 0,25).


В четвертой строке таблицы 7 рассматривается то же самое распределение


$\sigma$, однако здесь $N_1=N_2$, что приводит к симметричности кривой


$z(\varphi)$ относительно значения $\varphi=\pi/4$ (в частности, здесь


$\gamma_{10} = \gamma_{20}$). В последних двух строках таблицы 7


матрицы $\sigma$ отличаются перестановкой двух строк.


Видно, что опорные точки при этом остаются неизменными, а сами границы


несколько меняются. Максимальное различие в отклонении от базисной прямой


составляет для этих вариантов примерно 0,02.


\medskip





\hspace{104mm} Таблица 1


\smallskip





\begin{center}


\begin{tabular}{|c|c|c|c|c|c|c|}


\hline


$t$ & 0 & 0.1 & 0.2 & 0.3 & 0.4 & 0.45 \\


\hline


$\gamma_{10}$ & 1.49 & 2.02 & 3.00 & 5.30 & 13.84 & 33.44 \\


\hline


$\gamma_{20}$ & 1.45 & 1.93 & 2.79 & 4.70 & 11.14 & 24.96 \\


\hline


$z$ & $-0.21$ & $-0.33$ & $-0.59$ & $-1.25$ & $-4.04$ & $-10.64$ \\


\hline


\end{tabular}


\medskip


\end{center}





\hspace{119mm} Таблица 2


\smallskip





\begin{center}


\begin{tabular}{|c|c|c|c|c|c|c|}\hline


$\sigma$ &


\begin{tabular}{c}


1 0 0 0 \\ 1 0 0 0 \\ 1 0 0 0 \\


\end{tabular} &


\begin{tabular}{c}


1 1 0 0 \\ 1 1 0 0 \\ 1 0 0 0 \\


\end{tabular} &


\begin{tabular}{c}


1 1 0 0 \\ 1 1 0 0 \\ 1 1 0 0 \\


\end{tabular} &


\begin{tabular}{c}


1 1 1 0 \\ 1 1 1 0 \\ 1 1 0 0 \\


\end{tabular} &


\begin{tabular}{c}


1 1 1 0 \\ 1 1 1 0 \\ 1 1 1 0 \\


\end{tabular} &


\begin{tabular}{c}


1 1 1 1 \\ 1 1 1 0 \\ 1 1 1 0 \\


\end{tabular} \\


\hline


$\gamma_{10}$ & 0.593 & 0.593 & 0.690 & 0.690 & 1.171 & 1.171 \\


\hline


$\gamma_{20}$ & 0.586 & 0.586 & 0.586 & 0.586 & 0.586 & 0.689 \\


\hline


$z$ & $-0.00142$ & 0.03192 & $-0.00439$ & 0.07176 & $-0.03114$ & $-0.03817$\\


\hline


\end{tabular}


\medskip


\end{center}


\newpage





\hspace{117mm} Таблица 3


\smallskip





\begin{center}


\begin{tabular}{|c|c|c|c|} \hline


$\sigma$ & $\gamma_{10}$ & $\gamma_{20}$ & $z$ \\ \hline


\begin{tabular}{c}


1 1 1 1 1 1 1 1 0 \\ 1 1 1 1 1 1 1 1 0 \\ 1 1 1 1 1 1 1 1 0 \\


1 1 1 1 1 1 1 1 0 \\ 1 1 1 1 1 1 1 1 0 \\ 1 1 1 1 1 1 1 1 0 \\


1 1 1 1 1 1 1 1 0 \\


\end{tabular} &


\begin{tabular}{c}


1.1715729 \\ (1.1714449) \\


\end{tabular} &


\begin{tabular}{c}


0.5197831 \\ (0.5857864) \\


\end{tabular} &


\begin{tabular}{c}


$-0.0308206$ \\ ($-0.0311417$) \\


\end{tabular} \\


\hline


\end{tabular}


\medskip


\end{center}





\hspace{100mm} Таблица 4


\smallskip





\begin{center}


\begin{tabular}{|c|c|c|c|}


\hline


$\sigma$ & $\gamma_{10}$ & $\gamma_{20}$ & $z$ \\


\hline


\begin{tabular}{cccc}


1 & 1 & 1 & 1 \\ 1 & 0.45 & 1 & 1 \\ 1 & 1 & 1 & 1 \\


\end{tabular} & 22.1956 & 19.0499 & $-4.40151$ \\


\hline


\begin{tabular}{cccc}


1 & 1 & 1 & 1 \\ 1 & 0.49 & 1 & 1 \\ 1 & 1 & 1 & 1 \\


\end{tabular} & 118.09 & 99.01 & $-23.368$ \\


\hline


\end{tabular}


\medskip


\end{center}





\hspace{88mm} Таблица 5


\smallskip





\begin{center}


\begin{tabular}{|c|c|c|c|}


\hline


$\sigma$ & $\gamma_{10}$ & $\gamma_{20}$ & $z$ \\


\hline


\begin{tabular}{c}


1 0 1 0 1 \\ 1 0 1 0 1 \\ 0 1 0 1 0 \\ 1 0 1 0 1 \\


\end{tabular} & 1.000 & 0.685 & 0.056 \\


\hline


\end{tabular}


\medskip


\end{center}





\hspace{133mm} Таблица 6


\smallskip





%\begin{center}


\begin{tabular}{|p{2.2cm}|p{2.2cm}|p{2.2cm}|p{2.2cm}|p{2.2cm}|p{2.2cm}|}


\hline


$k$ & 1 & 2 & 3 & 4 & 5 \\


\hline


$z(\varphi(k))$ & $-7.9E-04$ & $-2.2E-04$ & $1.2E-003$ & $3.4E-003$


& $6.6E-003$ \\


\hline


\end{tabular}





\begin{tabular}{|p{2.2cm}|p{2.2cm}|p{2.2cm}|p{2.2cm}|p{2.2cm}|}


\hline


$k$ & 6 & 7 & 8 & 9 \\


\hline


$z(\varphi(k))$ & $1.2E-002$ & $2.0E-002$ & $3.5E-002$ & $5.6E-002$\\


\hline


\end{tabular}


\medskip


%\end{center}





\hspace{122mm} Таблица 7


\smallskip





\begin{center}


\begin{tabular}{|c|c|c|c|c|} \hline


$\sigma$ & $\gamma_{10}$ & $\gamma_{20}$ & $z$ \\


\hline


% \begin{tabular}{c} % if return, remove &


$\begin{smallmatrix}


\ & \ & \ & \ \\


0 &0 &0 &0 \\ 0 &0 &0 &0 \\ 0 &0 &0 &0 \\


\ & \ & \ & \ 


\end{smallmatrix}$ & 1.10557 & 1.17157 & $-3.33 E-008$ \\


\hline


% \begin{tabular}{c}


$\begin{smallmatrix}


\ & \ & \ & \  \\


0.24 &1.00 &0.24 &1.00 \\ 1.00 &0.24 &1.00 &0.24 \\


0.24 &1.00 &0.24 &1.00 \\


\ & \ & \ & \ 


\end{smallmatrix}$ & 55.2091 & 50.0000 & $-2.90177417$ \\


\hline


% \begin{tabular}{c}


$\begin{smallmatrix}


\ & \ & \ & \   \\


0.24 &1.00 &0.24 &1.00 \\ 1.00 &0.24 &1.00 &0.24 \\


0.24 &1.00 &0.24 &1.00 \\ 1.00 &0.24 &1.00 &0.24 \\


\ & \ & \ & \   


\end{smallmatrix}$ & 55.209068666 & 55.209068666 & $-4.08639082$ \\


\hline


% \begin{tabular}{c}


$\begin{smallmatrix}


\ & \ & \ & \ & \ & \ & \ & \ \\


0 &1 &1 &1 &1 &1 &1 &1 \\ 1 &0 &1 &1 &1 &1 &1 &1 \\ 1 &1 &0 &1 &1 &1 &1 &1 \\


1 &1 &1 &0 &1 &1 &1 &1 \\ 1& 1 &1 &1 &0 &1 &1 &1 \\


\ & \ & \ & \ & \ & \ & \ & \ 


\end{smallmatrix}$ & 2.666620894 & 2.663746419591 & $-4.6221435354E-001$ \\


\hline


% \begin{tabular}{c}


$\begin{smallmatrix}


\ & \ & \ & \ & \ & \ & \ & \ \\


0 &1 &1 &1 &1 &1 &1 &1 \\ 1 &0 &1 &1 &1 &1 &1 &1 \\ 1 &1 &0 &1 &1 &1 &1 &1\\


1 &1 &1 &1 &0 &1 &1 &1 \\ 1 &1 &1 &0 &1 &1 &1 &1 \\


\ & \ & \ & \ & \ & \ & \ & \ 


\end{smallmatrix}$ & 2.666620894 & 2.663746419591 & $-4.42714729620E-001$\\


\hline


\end{tabular}


\end{center}





\begin{thebibliography}{9}





\bibitem{1}


Самарский А.А. {\it Теория разностных схем.} -- 3-е


изд. -- М.: Наука, 1989. -- 616 с.


\bibitem{2}


Самарский А.А., Гулин А.В. {\it Устойчивость


разностных схем.} -- М.: Наука, 1973. -- 415 с.


\bibitem{3}


Гулин А.В., Самарский А.А. {\it Об устойчивости одного класса


разностных схем}


// Дифференц. уравнения. -- 1993. -- Т.\,29. -- \No\,7. -- С.\,1163--1174.


\bibitem{4}


Гулин А.В., Дегтярев С.Л. {\it Критерий устойчивости двуслойных и


трехслойных разностных схем}


// Дифференц. уравнения. -- 1996. -- Т.\,32. -- \No\,7. -- С.\,1--8.


\bibitem{5}


Самарский А.А., Гулин А.В., Вукославчевич В.


{\it Критерии  устойчивости двуслойных и трехслойных


разностных схем}


// Дифференц. уравнения. -- 1998. -- Т.\,34. -- \No\,7. -- С.\,975--979.


\bibitem{6}


Самарский А.А., Вабищевич П.Н., Матус П.П.


{\it Разностные схемы с операторными множителями.}


-- Минск: ЗАО ``ЦОТЖ'', 1998. -- 442 с.


\bibitem{7}


Гулин А.В., Юхно Л.Ф. {\it Численное исследование устойчивости


двуслойных разностных схем для двумерного уравнения теплопроводности}


// Журн. вычисл. матем. и матем. физ. -- 1996. -- Т.\,36.


-- \No\,8. -- С.\,118--126.


\bibitem{8}


Гулин А.В., Юхно Л.Ф. {\it Границы устойчивости двумерных


разностных схем}


// Матем. моделир. -- 1998. -- Т.\,10. -- \No\,1. -- С.\,44--50.


\bibitem{9}


Гулин А.В., Гулин В.А. {\it Границы устойчивости разностных схем с


переменными весовыми множителями}


// Изв. вузов. Математика. -- 1994. -- \No\,9. -- С.\,28--38.





\end{thebibliography}





\vskip 2cm





\post{\it Московский государственный университет}{\it Поступила}


\post{}{24.02.2000}





\label{end}


\end{document}


